\chapter{Can Human Rights Cross Cultures?}
\section{A Card That Cannot Speak for You}
Pick up something you may not yet own but have surely seen: an identity card or passport.

Hold it and look. It carries your name, birth date, photograph, and a number. This thin card can do nothing: take no beating for you, bear no hunger, leap to no defense when you are bullied. Yet it is enormously valuable. In some places, having it determines whether you may attend school, receive medical care, legally board a train, or be treated by police as a citizen rather than a problem.

A strange agreement hides inside. One side concerns the state, which issues the number, checks the photograph, and acknowledges your existence. The other concerns something less definite: how you deserve to be treated simply as a human being, without arbitrary confinement, beatings, or removal of your name. Oddly, this second agreement appears nowhere on the card, yet everyone assumes it is there.

Where did that agreement originate? Who may write it? Does its validity in China extend across the world? If local custom says a certain kind of person deserves particular treatment, does the agreement follow custom or stand against it?

This chapter concerns that unwritten second agreement, called human rights, and a debate continuing from the twentieth century: can human rights cross cultures?

\section{Paris, December 10, 1948}
Enter a scene with me.

Time: late afternoon, December 10, 1948. Place: the Palais de Chaillot beside the Seine in Paris. The brightly lit UN General Assembly hall, seating over a thousand, is full of delegates. Outside is genuine winter, but the preceding winter feels colder to most present. The Second World War ended only three years earlier. European rubble remains, and many among the tens of millions killed in camps, battles, and bombings still lack names on gravestones.

One item dominates today's agenda: a vote on the Universal Declaration of Human Rights.

It did not fall from the sky. A drafting committee has argued for nearly two years. Its chair, former US president's widow Eleanor Roosevelt, has presided with her gavel over dozens of overnight meetings. Lebanese philosopher-diplomat Charles Malik records and organizes text. China's Zhang Pengchun, a Columbia-educated educator, repeatedly breaks deadlocks by invoking Confucian ideas such as benevolence and not imposing on others what one does not want oneself. India's Hansa Mehta objects to the opening all men are born free and equal, demanding all human beings because men can exclude women. She wins. Today's first article uses the wording she secured.

The hall is not filled exclusively with European and American suits. Latin American delegations occupy a substantial area. Many come from countries recently escaping or still escaping dictatorship, knowing acutely what paper rights do and do not mean.

Voting begins: names called, hands raised. Forty-eight in favor, none against, then eight abstentions, including the Soviet Union and several allies, Saudi Arabia, and apartheid South Africa.

No opposing votes, eight abstentions. Remember this strange result. Nobody dares openly oppose human rights, yet eight hands stop halfway.

\section{Written Agreement, Unwritten Disagreement}
State the conflict before answering.

Outwardly, that day in 1948 marked a great reconciliation. Countries recently killing one another on a vast scale agreed on a list. Wherever someone is born, whatever their faith or language, certain basics cannot be taken away: life, liberty, freedom from slavery and torture, recognition as a person before law.

But eight abstentions pointed like fingers toward cracks beneath the table.

The Soviet bloc's argument was broadly that the list emphasized individual rights, free expression, and property, while someone without work, bread, or shelter would find freedom worth little. They compressed the disagreement into bread before ballots.

Saudi representatives did not openly reject the entire document at the vote, but certain provisions deeply troubled them: equal marriage rights for women and men conflicted with their interpretation of Islamic law, while freedom to change religion concerned not an individual right but a much graver matter of family and religious law.

South Africa's abstention was starkest. A country treating Black people as second-class citizens could hardly endorse universal equality.

Notice three distinct sources of unease: political philosophy about individual freedom and bread, religious tradition placing divine law above human lists, and naked vested interest threatened by equality. Calling them all opposition to human rights is too coarse. Here emerges the real question:

\textbf{What entitles a list naming everyone to govern groups that disagree?}

Those answering that it can face a sharp challenge: is the list humanity's consensus or the war's winners attaching a universal label to their values? Those answering that it cannot face an equally sharp challenge: where may a culture's oppressed people seek help---widows sacrificed at husbands' funerals, segregated Black people, daughters disposed of by families? From the very culture oppressing them?

Neither side answers easily, which makes the question worth a chapter. Before continuing, weigh whether universal equality is a discovered human fact or an idea some people invented and promoted worldwide.

\section[Abstentions: Two Full Arguments]{Behind the Abstentions: Two Full Arguments}
Present both positions at full strength, starting with the side often cast as villain. This deserves care: online discussion frequently equates cultural relativism with excusing oppression, unfairly.

\textbf{Position one: no court stands above culture.}

Make relativism's case fully. In 1947, a year before the vote, the American Anthropological Association issued a formal statement expressing concern about the draft. Anthropologists had spent much of the early twentieth century living in others' villages, learning their languages, and eating their food. Their deepest lesson was that a group's way of life forms an interconnected whole. You cannot remove one component and grade it backward by another system's standards. Their statement broadly argued that values grow from particular histories; one culture's answer to how people should live may not apply elsewhere. Calling a European-derived list universal might be another cultural arrogance, whose previous name was colonialism.

Three points give this strength. First, observation supports human difference: some societies prioritize family over individuals, afterlife over present life, or elders' authority. These are not merely uncivilized delays but complete arrangements operating for centuries. Second, relativism has a legitimate moral motive: it originally defended peoples Europeans called savages against being measured by outsiders' rulers. Third, it identifies a real danger: whoever interprets universal standards gains a useful weapon. Gunboats have often followed promises to spread civilization.

\textbf{Position two: oppressed people cannot wait.}

Universalism has a complete case too, likewise beginning with the weak. Notice that the strongest UN advocates were not exclusively Western powers, but often Latin American, Asian, and African representatives and activists from small countries. With authoritarian governments and compromised courts at home, domestic remedies were blocked. Standards above national governments offered not arrogance but a lifesaving ladder. They enabled an appeal to the world: this is not merely our internal affair or culture, but violation of an agreement our government itself accepted.

This side also has three arguments. First, pain crosses cultures: heat, hunger, and fear of arbitrary imprisonment need no translation. Definitions of good lives vary enormously, but torture, massacre, and starvation show striking commonality. A shared floor may be easier than a shared ceiling. Second, culture is not a solid block. Those invoking tradition are often its powerful members; women, children, poor people, and minorities were never asked. Cultural relativism can readily become a gag for internal victims. Third, the postwar world had seen how far regimes could stretch legality: concentration camps operated under law. Without standards above national laws, humanity would have no place from which even to say such regimes were wrong.

Widen the view and the dispute becomes more than West against non-West.

After the Cold War, several Asian governments promoted Asian values in the 1990s: family, order, collective interests, and development justified postponing individual political rights as another mature model, not a defect. In 1993, Asian states meeting in Bangkok emphasized historical and cultural contexts, development rights, and opposition to foreign interference under human-rights banners. Yet on the same continent, Myanmar's Aung San Suu Kyi publicly countered, in paraphrase: do not defend Asian prisons in Asian culture's name. Asian traditions also respect dignity; Burmese people are not a new species needing no freedom. That year, more than a hundred states at Vienna's World Conference on Human Rights reaffirmed universality and indivisibility. Most developing countries joined that affirmation. Universal versus relative is not an East-West map but a tug-of-war inside every country.

Revisit an easily mistaken assumption: the declaration as a Western document addressed to the world. In Paris, an Indian changed all men, a Chinese delegate introduced Confucian ideas, a Lebanese coordinated drafting, and dozens of affirmative hands came from Asia, Africa, and Latin America. Its sharpest academic critique, the anthropologists' statement, came from within the West. The picture is messier than Western promotion and non-Western resistance, and truth usually lies in that messier picture.

\section{Thinking Bridge: Four Questions About a Right}
Can we analyze claims of universal rights without allegiance or raised voices? Yes. Break every rights claim into four questions, skipping none.

\textbf{First: whose right?} Is the holder an individual or a group? Religious freedom usually concerns individuals; freedom of a people from extermination protects a group. Many arguments confuse these levels. Family honor treats the family as rights-holder above its members. Naming the holder can cool half the anger.

\textbf{Second: claimed against whom?} Rights do not float freely; each corresponds to someone's duty. Freedom from torture primarily addresses the state, which must neither torture nor permit others to do so. Education rights address governments and parents. Respect addresses everyone around us. Identifying the addressee reveals many conflicts as \textbf{duty-bearers passing responsibility}: government blames parents, parents schools, schools society.

\textbf{Third: on what grounds?} Distinguish legal grounds, meaning a law provides it; moral grounds, why it is right; and traditional grounds, what ancestors believed. A right may have all three or only some. Freedom from slavery now has all three; a weekly rest day once had only legal support in many places. Asking grounds blocks the fallacy that absence from written law means someone does not deserve a right.

\textbf{Fourth: who enforces it?} When violated, who acts: domestic courts, international tribunals, foreign armies, or public opinion alone? This painful question often disappoints: many rights are written, few enforceable. Yet do not rush into cynicism. Weak enforcement does not erase a right, any more than an undefended door proves doors should not exist. Rights are unfinished projects, not merely delivered inventories.

Before forwarding or attacking headlines about a country's abuses or a culture's disregard for rights, ask whose right, against whom, on what grounds, and enforced by whom. Heated disputes often reveal speakers inhabiting different questions.

\section{Reading Evidence: Articles 1 and 29}
Now read the text voted on December 10, 1948.

Article 1 of the Universal Declaration of Human Rights, adopted by the UN in 1948, reads:

\begin{quote}
All human beings are born free and equal in dignity and rights. They are endowed with reason and conscience and should act towards one another in a spirit of brotherhood.
\end{quote}
Read slowly: every word was contested. All includes every ethnicity, gender, faith, and fortune; Mehta fought over that scope. Born means rights are not state-issued awards for good behavior. Government recognizes rather than grants them, and therefore cannot legitimately withdraw them. Dignity precedes rights, an intriguing order: rights look like dignity itemized, a legal translation of treating people as people. According to participants' recollections, reason and conscience entered through contributions including Zhang Pengchun's. They invoke neither one religion's deity nor one school's doctrine, but a broadly recognizable idea: humans think and feel compassion.

Read another, rarely quoted but perhaps this chapter's most important: Article 29's opening, from the same source.

\begin{quote}
Everyone has duties to the community in which alone the free and full development of his personality is possible.
\end{quote}
Why duties in a declaration of rights? This is a crease left by negotiation. Rights are not a deposit account allowing withdrawals without deposits. You can speak, think, love, and hate because you grew within society, so claiming rights also entails owing society something. Later critics calling the declaration excessively individualistic often read only Article 1, not Article 29. The reminder is simple: \textbf{read a document completely before criticizing it}, a surprisingly scarce political habit.

Another frequently compared passage predates the declaration by more than two millennia. In the Analects' ``Yan Yuan,'' Zhonggong asks about benevolence, and Confucius replies:

\begin{quote}
Do not impose on others what you do not want for yourself.
\end{quote}
This is not equivalent to modern human rights: it concerns personal cultivation, not rights against the state. But reciprocal empathy as a restraint arose independently in distant traditions. Jewish and Christian teachings contain treating others as you wish to be treated; Islamic tradition has similar formulations. This is the Golden Rule. Its prevalence \textbf{does not prove} universal rights, but shows that justifying protection from arbitrary harm is no single culture's monopoly.

\section{Why It Emerged: Three Explanations}
Return to 1948. We know \textbf{what happened}: forty-eight affirmative votes, none against, eight abstentions, a declaration adopted, followed over ensuing decades by conventions against genocide and torture, covenants on civil-political and economic-cultural rights, and supervisory institutions. We broadly know \textbf{participants' understandings}: a new world's foundation to some, great-power packaging to others. Compare \textbf{why it happened}. Why did humanity create an everyone list in the mid-twentieth century? The following are competing explanations, not facts themselves.

\textbf{Explanation one: humanity learned a lesson, a moral-learning account.}

Wartime atrocities, especially industrialized extermination of a people, demonstrated evil's possible scale. People had thought civilization's progress would automatically restrain barbarism; now they knew otherwise. After painful reflection, the world wrote prohibited treatment into international law as an advance warning line. Many contemporary statements support this account. Nuremberg first brought crimes against humanity into legal practice, meaning some deeds remain criminal despite national permission or superior orders. Its limit is explaining half-learned lessons. Massacres continued, and great powers often ignored allies' atrocities. Learning was genuine, but authority to remind others was never equally distributed.

\textbf{Explanation two: power needed new clothing, a power-politics account.}

Ask who benefits rather than what is said. Postwar America and the Soviet Union competed for supremacy and brighter banners. Freedom and human rights became effective Cold War weapons, cheaply attacking the opposing camp. This explains what moral learning struggles with: weak enforcement because powers never intended treaties to bite themselves; accusations spotlighting opponents rather than allies; protection becoming an invitation to military intervention. Its limit is explaining why powerless people, small states, and domestic opposition desperately grasped the same rope. Once made, a tool does not belong solely to its maker. A blacksmith forges a blade, but anyone can hold it.

\textbf{Explanation three: the weak forced the door open, a social-pressure account.}

Move from great-power rooms to people outside. The UN Charter initially gave limited attention to rights. Dozens of smaller Latin American, Asian, and African states and civil groups---churches, unions, women's organizations, lawyers' associations---lobbied, pressured, and wrote memoranda to force provisions in. Over subsequent decades, rights gained teeth not through great-power generosity but through people struggling outside prisons, in courts, and on streets: doctors exposing torture, mothers listing the disappeared, lawyers defending political prisoners. This restores a crucial fact: institutional history is not only great-power competition. Its limit is explaining why equally desperate appeals sometimes move the world and sometimes vanish. Weak voices often become institutions only when powerful states find it convenient. That does not belittle the weak; it recognizes who holds the hinges.

Each explanation holds part of the truth. Comparing them reveals that a declaration full of admirable words can simultaneously be sincere repentance, a shrewd weapon, and a ladder purchased with nameless people's lives. None cancels the others. Learning to hold them together is among political history's hardest and most worthwhile tasks.

\section{Deeper Challenge: Thin Limits, Thick Reasons}
Now introduce a weighty political-philosophy tool for locating a position between cultural difference and shared minimums without merely smoothing over disagreement.

Michael Walzer distinguishes \textbf{thick and thin morality}. Within cultures, morality is thick: rooted in history, religion, language, and ritual, richly answering who we are and how to live. Such accounts differ and neither can nor need become uniform. Yet in extreme situations---massacre, torture, treatment of people as livestock---different cultures utter nearly identical brief words: stop, this is murder, this is unjust. These are thin, without elaborate doctrine or ritual, translating almost intact. Walzer observes that \textbf{humanity often shares a few prohibitions and minimums, not a complete thick value system}.

This loosens many knots. Asking everyone to accept one good life is impossible and undesirable; thick morality belongs to cultures' growth. Thin prohibitions of massacre, torture, and slavery need no single thick culture, with supporting resources available within any tradition. Thus more than a hundred states could affirm universality at Vienna in 1993: they endorsed minimums, not identical lifestyles.

Consider a frequently misused claim: a culture lacking the concept of rights cannot have human rights applied to it. Indian economist-philosopher and Nobel laureate Amartya Sen broadly counters in two ways. First, great traditions never speak with one voice. India has despotic emperors but also Ashoka, who ordered tolerance inscribed on stone. China has Legalism but also Mencius's ranking of the people first, the state next, and the ruler last, in ``Jin Xin II.'' Claims about a culture's essence often select its most convenient half. Second, even absence of the word rights would not disqualify people from protection. \textbf{The need for protection depends not on which word their language first contains, but on their capacity to hurt.}

Look farther. In many sub-Saharan African settings, ubuntu broadly means I am because we are: relationships and community fulfill personhood. It initially seems far from individual-rights lists. Yet after apartheid, this tradition's emphasis on relationship and repair supported South Africa's Truth and Reconciliation Commission, enabling victims to speak and society to confront their experiences. The thin prohibition of arbitrary harm remained, while its realization carried local thickness. Thin limits and thick reasons are not alternatives but foundation and building.

The theory has boundaries. Thin does not mean light: enforcing even a few minimums may demand enormous sacrifice. And who decides whether a norm is a minimum or local preference? Power quietly returns. Theory advances the difficulty without abolishing it. Be suspicious of any theory claiming to settle cultural difference completely.

\section[School Rules, Screens, and Distant War]{Today's Echoes: School Rules, Screens, and Distant War}
That vote over seventy years ago echoes nearby.

Schools stage miniature rights debates daily. May grades remain private? Does confiscating phones violate private space? Do cleaning duties, group exercise, and uniform hairstyles protect collective interests or diminish individuals? The four questions clarify many disputes: whose rights, against school or classmates, grounded in rules, law, or tradition, enforced by teachers, parents, or nobody? Hardest are conflicts with genuine values on both sides: one person's quiet study and a whole class's quiet study. This is the smallest model of the individual-community-power triangle, present in classrooms without visiting the UN.

On screens, rights simultaneously expand and contract. Videos of distant beatings, wrongful imprisonment, and forced demolitions reach hundreds of millions within minutes, weakening nobody-knows as a cover. Meanwhile, enormous quantities of locations, conversations, and browsing preferences are collected, with privacy clauses scarcely read before agreement. Selection is subtler: one war reaches your home page, another equally large conflict may never appear in a year. Recall the explanations: human-rights spotlights follow not only darkness but whoever holds the lamp. Platforms now hold part of it.

Farther away, the 1994 genocide against Tutsi in Rwanda killed an often-cited estimated eight hundred thousand people in around a hundred days, with little international prevention despite warnings. In 1995, around eight thousand Muslim men and boys were killed in Srebrenica, Bosnia, within a UN-designated safe area. These are painful footnotes to minimum consensus: worldwide agreement against genocide could not make ink stop blades. But the other half matters. Such failures helped generate the idea that sovereignty carries responsibility, not unlimited permission behind closed borders. When a state massacres its people, protection's responsibility extends outward. This remains fiercely disputed: supporters see minimums gaining teeth, critics a potential great-power intervention license. Both concerns are genuine. Mechanisms can protect people or selectively target rivals. No perfect safeguard against selection exists, only an imperfect development: more people compare cases and ask why intervention happens here but not there.

Another nearby scene is dinner. An elder says children cannot discuss rights with adults because these are family rules. How might you respond? Reciting the declaration probably fails and reinforces claims that rights are foreign. Perhaps ask whom family rules protect and constrain, and where the weakest member can seek help when hurt most. Once the question becomes who may call on whom, culture's great wall lowers, revealing particular people behind it.

\section{For You: Who Draws the Minimum Line?}
Return to the card.

Identity cards and passports confirm state recognition. Human rights claim that even without any state's recognition---as a refugee, stateless person, or someone stripped of nationality---you deserve human treatment. If this agreement comes from no passport, where does it come from, and who sets its boundaries?

Before answering, weigh the chapter's evidence.

Relativism offers real cultural histories, unequal authority to interpret universal standards, and their repeated use as conquerors' clothing. Making one lifestyle humanity's destination is itself violence.

Universalism offers pain and humiliation needing no translation, powerful people speaking for culture while victims cannot appeal, and thin minimums---no massacre, torture, or treatment as livestock---supported in almost every tradition and affirmed by diverse representatives in 1948 and 1993.

You also know harder facts: enforcement perpetually trails written promises; protection can be selectively used by powerful states, yet returning everything to national governments removes the weak person's ladder. Development before freedom sometimes describes circumstances and sometimes postpones freedom indefinitely.

Is the claim that being human entitles someone to certain things a fact humanity gradually discovers, or a shared vow made at a terrible historical moment? If discovered, why was it buried so long? If a vow, why does it bind those absent when sworn?

Finally, when you say this is my right, do you recognize you also assign another person a duty? Will you shoulder an equal duty yourself?

There is no standard answer. But your response shapes whom you become tomorrow in classrooms, group chats, and encounters with news of distant war.
